\documentclass{article}
\usepackage{graphicx}
\usepackage{amsmath}
\usepackage{amssymb}
\usepackage{amsfonts}
\usepackage{textcomp}
\usepackage{amsthm}
\usepackage{mathtools}

\theoremstyle{plain}
\newtheorem{theorem}{Theorem}[section]
\newtheorem{lemma}[theorem]{Lemma}
\newtheorem{corollary}[theorem]{Corollary}
\newtheorem{proposition}[theorem]{Proposition}

\theoremstyle{definition}
\newtheorem{definition}[theorem]{Definition}
\newtheorem{example}[theorem]{Example}

\theoremstyle{remark}
\newtheorem{remark}[theorem]{Remark}

\title{P-03}
\author{A. E. Mahrouss\\amlal@nekernel.org}
\date{\today}

\begin{document}
\maketitle

\section{Theorems}

\begin{theorem}
Suppose
\begin{equation}
    A(n, p) \coloneq \sum_{n}^{p} M(n p)
\end{equation}
\begin{equation}
    A(n, p) \implies p! + \text{ } (n+1p)! + \text{ } (n+2p)! + \text{ } \dots \text{ } + (n+Lp)! 
\end{equation}
For $A(n, p) \neq 0$, the previous identity holds for numbers in $\mathbb{R}$.
\end{theorem}

\end{document}